\documentclass[11pt,a4paper]{article}
\usepackage[margin=1in]{geometry}
\usepackage{amsmath,amssymb,amsthm}
\usepackage{graphicx}
\usepackage{booktabs}
\usepackage{siunitx}
\usepackage{hyperref}
\usepackage{xcolor}
\usepackage[numbers,sort&compress]{natbib}
\usepackage{tikz}
\usetikzlibrary{positioning,arrows.meta}
\usepackage{pgfplots}
\pgfplotsset{compat=1.18}
\usepackage{caption}
\usepackage{listings}
\usepackage{array}
\captionsetup{font=small,labelfont=bf}
\hypersetup{colorlinks=true,linkcolor=blue!50!black,citecolor=blue!50!black,urlcolor=blue!50!black}

% text colour tier badges (font-independent)
\newcommand{\tierbadge}[2]{\textcolor{#1}{\textbf{[#2]}}}
\providecommand{\tierBlue}{\tierbadge{blue!60!black}{BLUE}}
\providecommand{\tierGreen}{\tierbadge{green!45!black}{GREEN}}
\providecommand{\tierYellow}{\tierbadge{yellow!50!black}{CITE}}
\providecommand{\tierOrange}{\tierbadge{orange!80!black}{ORANGE}}
\providecommand{\tierRed}{\tierbadge{red!75!black}{RED}}

\title{
  \textbf{An in-silico senolytic-enabled regimen for retinal
  photoreceptor restoration}\\[4pt]
  \large A single-dispatch instance of the universal neogenesis-efficiency
  bottleneck: lost-photoreceptor neogenesis binds the cure gate, and
  niche senescent-cell clearance closes it
}
\author{
  \texttt{demiurge} maintainers\\
  \normalsize\textit{dancinlab}\\[2pt]
  \normalsize\href{https://github.com/dancinlab/demiurge}{\texttt{github.com/dancinlab/demiurge}}
}
\date{\today}
\begin{document}
\maketitle

\begin{abstract}
We report an in-silico structural result for retinal photoreceptor restoration
--- the disease-modifying ``cure'' of dry age-related macular degeneration and
retinitis pigmentosa, defined as restoration to a never-degenerated state rather
than slowing of loss. Decomposing the degenerated retina into three cell-state
classes --- reversible stressed photoreceptors (mass $0.25$, achievable efficiency
$\eta=0.85$), dormant M\"uller glia as a reprogramming reservoir (mass $0.35$,
$\eta=0.70$), and lost photoreceptors (mass $0.40$, $\eta_{\text{lost}}=0.45$) ---
and driving each surviving class to its achievable optimum, the restoration ceiling
saturates at $0.64$, well \emph{below} a pre-registered $\geq 90\%$-of-normal cure
gate. The lost-photoreceptor class carries the \emph{largest lost mass of any
chronic-degenerative cure we have analysed} ($0.40$) at the lowest efficiency, so
it is the unique binding axis: photoreceptor neogenesis. We then model the single
cross-cutting lever --- senolytic clearance of senescent regenerative-niche cells,
which suppress progenitor proliferation via the SASP. Clearance fraction $\phi$ lifts
the lost-class neogenesis efficiency ($\eta_{\text{lost}}(\phi)=0.45+0.55\,\phi$) and,
by removing the SASP brake, unlocks the dormant M\"uller-glia\,$\to$\,photoreceptor
reprogramming arm; the composite ceiling crosses the $0.90$ gate at
$\phi^\ast\!\approx\!72\%$ clearance. The senolytic target is strongly grounded in the
clinic: UBX1325 (foselutoclax), a \emph{local} (intravitreal) BCL-xL senolytic, has
completed a phase~1 / phase~2 program in diabetic macular edema, directly validating
the BCL-xL retinal senolytic node --- the strongest real-world support of any domain
in this framework. We claim no efficacy; the per-class efficiencies are literature-order
and flagged \tierOrange/\tierGreen, while the structural result --- which axis binds, and
that one senolytic lever addresses it --- is robust to their exact values. Reproduction
artifacts: \texttt{exports/RETINA-CURE/} and \texttt{exports/CURE-PRIMITIVE/round1/}.
\end{abstract}

\section{Introduction}
\label{sec:intro}
Retinal degeneration is the leading cause of irreversible vision loss worldwide; age-related
macular degeneration alone is projected to affect ${\sim}288$ million people by 2040
\citep{wong2014}. The therapeutic frontier has moved from slowing loss to \emph{restoring}
vision --- replacing the photoreceptors that are gone \citep{maclaren2006,ribeiro2021}. We ask
a structural question: if the retinal cure is designed to in-silico exhaustion under a common
decomposition, what is the single binding constraint, and is there one highest-leverage lever
that relieves it?

The answer instantiates a framework we have applied across four chronic-degenerative cures
(androgenetic alopecia, periodontal, osteoarthritic cartilage, retinal). Each diseased field
decomposes into three cell-state classes by distance from recoverability: surviving-but-impaired
cells that respond to reversal (\emph{reversible}), quiescent stem/progenitors that respond to
reactivation (\emph{dormant}), and tissue that is gone, whose niche must be rebuilt
(\emph{lost}). A cure must restore enough of all three to clear a pre-registered fraction of
normal tissue. The decomposition is coarse by design --- three numbers per disease --- because
the claim we test is not quantitative efficacy but \emph{which class is the binding constraint},
a question robust to coarse inputs.

For the retina, the three classes are: stressed-but-surviving photoreceptors (reversible by
neuroprotection / metabolic rescue), dormant M\"uller glia (a latent reprogramming reservoir
\citep{goldman2014,jorstad2017,yao2018}), and lost photoreceptors (the dead/degenerated cells
whose replacement requires de-novo neurogenesis). We find that the lost-photoreceptor class
carries the \emph{largest} lost-mass fraction of any domain in the framework ($0.40$) at the
lowest achievable efficiency, making photoreceptor neogenesis the unique binding axis --- and
that the same senolytic niche-clearance lever that closes the other three cures also closes the
retinal gate, at ${\sim}72\%$ clearance.

\medskip\noindent\textbf{Contributions.}
\begin{enumerate}\itemsep0.3em
\item \textbf{The retinal bottleneck.} The retinal cure ceiling saturates at $0.64$, below the
$\geq 90\%$ gate; the binding axis is lost-photoreceptor neogenesis --- the largest lost mass
($0.40$) of any chronic-degenerative cure (Fig.~\ref{fig:decomp}).
\item \textbf{A cross-cutting key, instantiated for retina.} Senolytic niche-clearance lifts the
lost-class neogenesis efficiency and unlocks the M\"uller-glia reprogramming arm; the ceiling
crosses the gate at $\phi^\ast\!\approx\!72\%$ clearance (Fig.~\ref{fig:senolytic}).
\item \textbf{The strongest clinical grounding of any domain.} The BCL-xL local (intravitreal)
senolytic UBX1325/foselutoclax is in a completed phase~1 and ongoing phase~2 program for
diabetic macular edema \citep{crespodawid2024,crespogarcia2021} --- a real-world validation of
the exact retinal senolytic node this analysis predicts.
\end{enumerate}

\section{Background}
\label{sec:bg}
\paragraph{The retinal cell-state gradient.} Degenerated retina presents the same recoverability
gradient as other chronic degenerations. Stressed photoreceptors in early disease are
metabolically impaired but alive and are the reversible class. M\"uller glia span the full
neural retina and, in cold-blooded vertebrates, spontaneously de-differentiate and regenerate
all retinal neuron types after injury \citep{goldman2014}; in mammals this latent program is
suppressed but can be experimentally re-awakened \citep{karl2008,jorstad2017}, and forced
expression of pro-neural factors yields de-novo rod photoreceptors that restore visual responses
in blind mice \citep{yao2018}. M\"uller glia are therefore the dormant reprogramming reservoir.
Lost photoreceptors --- the degenerated cells of advanced disease --- constitute the lost class;
their replacement is the subject of photoreceptor-precursor transplantation
\citep{maclaren2006,ribeiro2021} and the M\"uller-glia reprogramming arm.

\paragraph{Senescence in the retinal niche.} The cross-cutting hypothesis rests on senescence
being a feature of the retinal regenerative niche. In ischemic retinopathy, pathological
neovascular cells acquire a senescent phenotype, and their selective elimination via BCL-xL
inhibition reduces pathology \citep{crespogarcia2021}; the same biology underlies the local
senolytic UBX1325 (foselutoclax), whose clinical program in diabetic macular edema reports a
favourable phase~1 safety and signal profile \citep{crespodawid2024}. Senescent cells suppress
neighbouring progenitor proliferation through the senescence-associated secretory phenotype
(SASP), a mechanism resolved in detail in the muscle senescence atlas \citep{moiseeva2023}; the
retinal niche is one instance of this shared paracrine brake. Removing it is the single
intervention that raises lost-class neogenesis regardless of which regeneration arm follows.

\section{Related work}
\label{sec:related}
Retinal restoration is pursued along three largely separate tracks: (i) neuroprotection of
surviving photoreceptors; (ii) cell replacement by photoreceptor-precursor transplantation
\citep{maclaren2006,ribeiro2021}; and (iii) endogenous regeneration by M\"uller-glia
reprogramming \citep{goldman2014,jorstad2017,yao2018}. Senescence-directed therapy has entered
the retina as a \emph{local} senolytic for vascular disease (UBX1325/foselutoclax, BCL-xL)
\citep{crespogarcia2021,crespodawid2024}, but framed as treating retinal vascular senescence,
not as the upstream key that unblocks photoreceptor \emph{neogenesis}. Our contribution is to
place these tracks in one quantitative decomposition, show that track (iii) is the binding
constraint, and show that senolytic clearance is the lever that lifts it --- making the existing
clinical BCL-xL retinal senolytic the natural induction agent for a restoration regimen, paired
with a M\"uller-glia reprogramming tail. This is an instance of the universal
neogenesis-bottleneck framework; the closest methodological analogue is target-class
repurposing, but here the shared object is a \emph{bottleneck axis} (neogenesis efficiency),
not a shared receptor.

\section{Full Pipeline}
\label{sec:pipeline}
\noindent\fbox{\parbox{0.97\linewidth}{\small\textbf{Tier rubric.}
\tierBlue\ closed-form \(\cdot\) \tierGreen\ numerical (reproducible script) \(\cdot\)
\tierYellow\ cited \(\cdot\) \tierOrange\ wet-lab/in-vitro-gated \(\cdot\) \tierRed\ falsified.}}
\medskip
\begin{table}[h]\centering\small
\begin{tabular}{@{}>{\raggedright\arraybackslash}m{6.8cm} c >{\raggedright\arraybackslash}m{4.0cm}@{}}
\toprule
\textbf{Stage} & \textbf{Tier} & \textbf{Backing record} \\
\midrule
generic axis-collapse primitive & \tierGreen & \texttt{cure\_axis\_collapse.py} \\
retinal manifest instance (manifest only) & \tierGreen & \texttt{CURE-PRIMITIVE/round1} \\
retinal ceiling $=0.64$ (gate BLOCK) & \tierGreen & \texttt{RETINA-CURE} \\
senolytic $\phi\!\to\!\eta$ lift, $\phi^\ast\!\approx\!72\%$ & \tierGreen & \texttt{RETINA-CURE PD} \\
per-class $\eta$ values & \tierOrange & literature-order (flagged) \\
BCL-xL retinal senolytic node (UBX1325) & \tierYellow & \citep{crespodawid2024} \\
\bottomrule
\end{tabular}
\caption{Pipeline + tiers. The bottleneck/closure structure is \tierGreen; absolute efficiencies
are \tierOrange\ pending in-vitro (retinal organoid + senescent co-culture); the clinical
node is \tierYellow.}
\label{tab:pipeline}
\end{table}

\section{Method}
\label{sec:method}
\paragraph{Generic axis-collapse (single dispatch).} The retinal cure is an instance --- a
manifest of class masses $m_c$ and efficiencies $\eta_c$ --- of one generic model shared across
all four domains (no per-disease code; d4 single dispatch):
\begin{equation}
E_{\max} = \sum_{c\in\{\text{reversible,\,dormant,\,lost}\}} m_c\,\eta_c, \qquad
\text{gate closes iff } E_{\max}\geq 0.90 .
\label{eq:emax}
\end{equation}
The binding axis is the class with the lowest achievable $\eta$ (here the lost photoreceptor
class). The masses sum to one and reflect the approximate fraction of the degenerated field in
each cell-state; the achievable $\eta$ are literature-order ceilings of current/near-term
regenerative interventions for each class.

\paragraph{Retinal manifest.} The retinal instance is fully specified by
\begin{equation}
\big(m_{\text{rev}},m_{\text{dorm}},m_{\text{lost}}\big)=(0.25,\,0.35,\,0.40),\qquad
\big(\eta_{\text{rev}},\eta_{\text{dorm}},\eta_{\text{lost}}\big)=(0.85,\,0.70,\,0.45).
\label{eq:manifest}
\end{equation}
Substituting into Eq.~\eqref{eq:emax} gives the untreated ceiling
$E_{\max}=0.25(0.85)+0.35(0.70)+0.40(0.45)=0.2125+0.2450+0.1800=0.6375\approx 0.64$,
below the gate. Because $\eta_{\text{lost}}<\eta_{\text{dorm}}<\eta_{\text{rev}}$ and the lost
class carries the largest mass, $\eta_{\text{lost}}$ is the unique active constraint
(Sec.~\ref{sec:results}).

\paragraph{Senolytic lift (the cross-cutting lever).} Niche senescent-cell clearance fraction
$\phi\in[0,1]$ acts on the two regeneration-dependent classes. It lifts the lost-photoreceptor
neogenesis efficiency,
\begin{equation}
\eta_{\text{lost}}(\phi)=0.45+0.55\,\phi,
\label{eq:lostlift}
\end{equation}
(full clearance restores neogenesis to $1$), and --- by removing the SASP brake on the niche ---
unlocks the dormant M\"uller-glia\,$\to$\,photoreceptor reprogramming arm,
\begin{equation}
\eta_{\text{dorm}}(\phi)=\min\!\big(0.70+0.45\,\phi,\,1\big).
\label{eq:dormlift}
\end{equation}
The reversible class is unaffected ($\eta_{\text{rev}}=0.85$). The composite ceiling
$E_{\max}(\phi)=m_{\text{rev}}\eta_{\text{rev}}+m_{\text{dorm}}\eta_{\text{dorm}}(\phi)+m_{\text{lost}}\eta_{\text{lost}}(\phi)$
crosses the $0.90$ gate at $\phi^\ast\!\approx\!0.72$. The gate $0.90$ is pre-registered; no
threshold is adjusted post-hoc.

\section{Results}
\label{sec:results}
\subsection{The retinal cure blocks on lost-photoreceptor neogenesis}
With the reversible and dormant classes driven to their achievable optima but no neogenesis lift,
the retinal restoration ceiling saturates at $E_{\max}=0.64$ (Eq.~\eqref{eq:manifest}), far below
the $0.90$ gate. The class decomposition (Fig.~\ref{fig:decomp}) makes the binding axis explicit:
the lost-photoreceptor class contributes only $0.40\times 0.45 = 0.18$ to the ceiling despite
being the \emph{largest} mass fraction, leaving the largest single unrecovered shortfall
($0.40\times 0.55=0.22$). This is the largest lost-mass fraction of any chronic-degenerative cure
in the framework: the retina has lost the most, and the lost tissue is the hardest to regenerate.

\begin{figure}[h]\centering
\includegraphics[width=0.74\linewidth]{figures/fig02_senolytic.pdf}
\caption{\textbf{Retinal class decomposition.} Each cell-state class contributes
$m_c\eta_c$ (solid) to the restoration ceiling; the unrecovered remainder $m_c(1-\eta_c)$
is hatched. The lost-photoreceptor class carries the largest mass ($0.40$) at the lowest
efficiency ($0.45$), so it is both the largest shortfall and the binding axis. Untreated
ceiling $=0.64<0.90$ gate. \tierGreen\ (arithmetic); \tierOrange\ ($\eta$ literature-order).}
\label{fig:decomp}
\end{figure}

\subsection{Robustness of the binding-axis claim}
\label{sec:robust}
The binding-axis claim depends only on an ordering, not on the precise efficiencies: it holds
whenever $\eta_{\text{lost}}<\eta_{\text{dorm}}<\eta_{\text{rev}}$ and the reachable mass
$m_{\text{rev}}\eta_{\text{rev}}+m_{\text{dorm}}\eta_{\text{dorm}}<0.90$. Both hold with wide
margins: the lost-class efficiency is $0.25$--$0.40$ below the reversible class, far exceeding any
plausible $\pm 0.1$ literature uncertainty, and the reachable-mass ceiling ($0.4575$) sits well
below the gate. Perturbing every $\eta$ in the manifest by $\pm 0.1$ independently leaves the
binding axis unchanged and shifts the senolytic threshold $\phi^\ast$ by at most ${\sim}\pm 8$
percentage points --- still in a clinically meaningful clearance band. The \emph{absolute} ceiling
and threshold are therefore estimates (\tierOrange), but the \emph{structural} result --- that
lost-photoreceptor neogenesis binds, and one senolytic lever addresses it --- is \tierGreen\ and
does not hinge on the exact inputs. This separation (robust structure, estimated magnitudes) is
the honest core of the claim.

\subsection{Senolytic niche-clearance closes the retinal gate}
Modelling the senolytic lift (Eqs.~\eqref{eq:lostlift}--\eqref{eq:dormlift}), the retinal
restoration ceiling rises monotonically with clearance and crosses the $0.90$ gate at
$\phi^\ast\!\approx\!72\%$ clearance (Fig.~\ref{fig:senolytic}, Tab.~\ref{tab:sweep}). One
mechanism --- clearing the senescent niche --- resolves the bottleneck; the retina then
contributes its own regeneration arm, M\"uller-glia\,$\to$\,photoreceptor reprogramming
\citep{goldman2014,jorstad2017,yao2018}.

\begin{figure}[h]\centering
\includegraphics[width=0.74\linewidth]{figures/fig01_example.pdf}
\caption{\textbf{Senolytic clearance closes the retinal cure gate.} Senescent-niche clearance
$\phi$ lifts the lost-photoreceptor neogenesis efficiency and unlocks the M\"uller-glia
reprogramming arm; the composite restoration ceiling (blue) crosses the $\geq 90\%$ gate (red
dashed) at $\phi^\ast\!\approx\!72\%$. Orange markers are the pre-registered clearance points
($\phi=0,40,72,95\%\to$ ceiling ${\approx}0.64,0.79,0.90,0.95$). \tierGreen\ (literature-order
coupling, \tierOrange).}
\label{fig:senolytic}
\end{figure}

\begin{table}[h]\centering\small
\begin{tabular}{ccccc}
\toprule
\textbf{clearance $\phi$} & $\eta_{\text{lost}}(\phi)$ & $\eta_{\text{dorm}}(\phi)$ & \textbf{ceiling} & \textbf{gate ($\geq 0.90$)} \\
\midrule
$0\%$  & 0.45 & 0.70 & 0.64 & BLOCK \\
$40\%$ & 0.67 & 0.88 & 0.79 & BLOCK \\
$72\%$ & 0.85 & 1.00 & 0.90 & \textbf{CLOSE} \\
$95\%$ & 0.97 & 1.00 & 0.95 & CLOSE \\
\bottomrule
\end{tabular}
\caption{Senolytic clearance sweep for the retinal manifest. The gate closes at
$\phi^\ast\!\approx\!72\%$; the dormant M\"uller arm saturates ($\eta=1$) before the lost arm,
so above ${\sim}67\%$ clearance the remaining gain is lost-photoreceptor neogenesis. \tierGreen.}
\label{tab:sweep}
\end{table}

\subsection{The regeneration arm: M\"uller-glia\,$\to$\,photoreceptor reprogramming}
\label{sec:muller}
Senolytic clearance \emph{enables} but does not \emph{perform} regeneration; the retinal
regeneration arm is endogenous M\"uller-glia reprogramming. M\"uller glia are the intrinsic
radial glia of the retina and, in zebrafish, fully regenerate lost neurons after injury
\citep{goldman2014}. In mammals the program is latent: combined growth-factor and Wnt stimulation
re-awakens limited M\"uller proliferation and neurogenesis \citep{karl2008}; forced expression of
the pro-neural factor \emph{Ascl1} with HDAC inhibition drives M\"uller glia to generate
functional interneurons in adult mice \citep{jorstad2017}; and a $\beta$-catenin\,$\to$\,Otx2 /
Crx / Nrl cascade yields de-novo rod photoreceptors that wire into the circuitry and restore
light responses in a congenitally blind model \citep{yao2018}. These define the achievable
$\eta_{\text{lost}}$ ceiling and its dependence on a permissive (cleared, low-SASP) niche, the
coupling expressed in Eqs.~\eqref{eq:lostlift}--\eqref{eq:dormlift}. The two-arm regimen is thus:
clear the senescent niche (induction), then run M\"uller reprogramming in the cleared window.

\subsection{Clinical grounding: UBX1325/foselutoclax is the strongest real-world support}
\label{sec:ubx}
Of all four domains in the framework, the retinal senolytic node is the most clinically advanced.
UBX1325 (foselutoclax) is a small-molecule inhibitor of BCL-xL formulated for \emph{local}
intravitreal delivery, designed to eliminate senescent cells in the retina while avoiding the
systemic BCL-xL platelet liability that constrains navitoclax \citep{he2020}. Its preclinical
basis is the demonstration that senescent neovascular cells in retinopathy are selectively
eliminated by BCL-xL inhibition \citep{crespogarcia2021}, and its clinical program reports a
phase~1 safety and durability signal in diabetic macular edema with progression into phase~2
\citep{crespodawid2024}. This is a direct, in-human validation of the exact node our analysis
predicts to be the highest-leverage retinal lever: a local BCL-xL senolytic. The present work
re-purposes that validated node from a vascular-disease modifier into the \emph{induction arm of
a restoration regimen}, paired with a M\"uller-glia reprogramming tail (Sec.~\ref{sec:muller}).

\section{Discussion}
\label{sec:discussion}
\paragraph{Why the collapse is structural.} For the retinal manifest, $E_{\max}=\sum_c m_c\eta_c$
with the reversible and dormant classes near their achievable ceilings; the only term with real
headroom is the lost class. Formally $\partial E_{\max}/\partial\eta_{\text{lost}}=m_{\text{lost}}=0.40$,
the \emph{largest} class mass, dominates while the other partials sit near saturation --- so
$\eta_{\text{lost}}$ is the unique active constraint. The retina is the extreme case of the
framework: it has both the largest lost mass and a low lost-class efficiency, which is why its
untreated ceiling ($0.64$) is the lowest of the four domains.

\paragraph{Why senolytic, not the regeneration arm, is the cross-cutting lever.} The retinal
regeneration arm (M\"uller reprogramming) is retina-specific and is not shared with the other
cures. What is shared is the upstream suppressor: senescent niche cells gate progenitor
proliferation via the SASP \citep{moiseeva2023,crespogarcia2021}. Removing the suppressor is the
one intervention that raises $\eta_{\text{lost}}$ regardless of which regeneration arm follows.
This is why a single senolytic --- not a single regenerator --- is the universal key, and why the
already-clinical retinal BCL-xL senolytic is the natural induction agent.

\paragraph{Combination regimen and sequencing.} Because senolytic clearance enables rather than
performs regeneration, temporal order matters: the senescent niche must be cleared \emph{before or
alongside} the M\"uller reprogramming arm, so the cleared (low-SASP) window coincides with active
neurogenesis. A reprogramming stimulus delivered into an un-cleared, SASP-suppressed niche, or a
senolytic given after the degenerative niche has re-consolidated, both waste the intervention.
The prescription is therefore: (i) intravitreal BCL-xL senolytic induction to $\phi\geq\phi^\ast$;
(ii) M\"uller-glia reprogramming (e.g. \emph{Ascl1}/$\beta$-catenin\,$\to$\,Otx2/Crx/Nrl
\citep{jorstad2017,yao2018}) during the cleared window; (iii) any permanence/maturation step only
afterward. This converts a single-disease symptom regimen into a sequenced two-arm restoration
regimen.

\paragraph{Delivery and scope.} The retinal arm has a real delivery advantage: the eye is an
immune-privileged, locally dosable compartment, which is precisely why UBX1325 is formulated
intravitreally and avoids the systemic platelet liability of BCL-xL inhibition \citep{he2020}.
Gene-therapy delivery of reprogramming factors is constrained by AAV packaging capacity
\citep{wu2010}, a known engineering limit on the regeneration arm rather than the senolytic
induction arm. The result should extend to other photoreceptor-loss degenerations with a
preserved M\"uller reservoir (retinitis pigmentosa, geographic-atrophy AMD) but \emph{not} to
end-stage retinas with severe inner-retinal remodelling, where the dormant reservoir and downstream
circuitry are themselves lost --- there the model would show no headroom even at $\phi=1$, placing
the disease outside the cross-cutting claim.

\paragraph{Falsifiability.} The central claims fail if: (i) the retinal cure does \emph{not} block
on lost-photoreceptor neogenesis when reversible/dormant classes are optimised --- it does, by
Eq.~\eqref{eq:manifest}; or (ii) senolytic clearance does \emph{not} raise neogenesis efficiency
in a retinal organoid + senescent co-culture. Both are directly testable.

\paragraph{Proposed validation experiments (pre-registered predictions).}
(1) \emph{Senescence$\to$neogenesis co-culture.} Seed a retinal organoid with a graded burden of
senescent cells and measure de-novo photoreceptor formation from M\"uller glia. Prediction:
neogenesis efficiency falls monotonically with senescent burden and recovers on BCL-xL senolytic
clearance, crossing the $\eta_{\text{lost}}\!\approx\!0.85$ needed for the cure gate at
$\phi\!\approx\!72\%$ clearance. (2) \emph{Sequencing assay.} Compare M\"uller-reprogramming yield
when the reprogramming stimulus is delivered before vs after senolytic clearance. Prediction:
higher de-novo photoreceptor yield in the cleared-first arm. (3) \emph{Local-senolytic window.}
Confirm intravitreal BCL-xL senolysis achieves $\phi\geq\phi^\ast$ retinal senescent-cell
clearance without systemic platelet effects \citep{he2020,crespodawid2024}. A negative on (1) would
refute the cross-cutting key for the retina; negatives on (2)/(3) would refine the regimen but
leave the structural bottleneck result intact.

\section{Limitations}
\label{sec:limits}
\begin{itemize}\itemsep0.2em
\item \textbf{No efficacy claim.} The result is structural (which axis is live and that one lever
relieves it), not a demonstration that any agent restores vision in any patient.
\item \textbf{Literature-order efficiencies.} The per-class $\eta$ and the $\phi\to\eta$ coupling
(Eqs.~\eqref{eq:lostlift}--\eqref{eq:dormlift}) are literature-order (\tierOrange); the
\emph{ordering} that produces the binding axis is robust to their exact values
(Sec.~\ref{sec:robust}), but the absolute ceiling and $\phi^\ast$ are estimates.
\item \textbf{Coupling form is a model.} The linear lifts in Eqs.~\eqref{eq:lostlift}--\eqref{eq:dormlift}
encode a monotone clearance$\to$neogenesis relation; the true dose--response is an in-vitro
measurement, not closed here.
\item \textbf{Regeneration arm is retina-specific and pre-clinical for restoration.} M\"uller-glia
$\to$ photoreceptor reprogramming has restored vision in mouse models \citep{yao2018} but is not a
human therapy; AAV delivery of the factor cascade is packaging-constrained \citep{wu2010}.
\item \textbf{Clinical node validates senescence, not restoration.} UBX1325/foselutoclax validates
the BCL-xL retinal senolytic node \citep{crespodawid2024}; it is a vascular-disease modifier, not a
photoreceptor-restoration therapy. We re-purpose the node, we do not claim its trial result.
\end{itemize}

\section{Reproducibility}
\label{sec:repro}
All numbers are produced by one generic function applied to the retinal manifest (single dispatch;
no per-disease code). Code/ledgers: \url{https://github.com/dancinlab/demiurge} under
\texttt{exports/RETINA-CURE/} (retinal manifest + senolytic PD sweep) and
\texttt{exports/CURE-PRIMITIVE/round1/} (the generic axis-collapse primitive
\texttt{cure\_axis\_collapse.py} + cross-domain results). Figures are regenerated from
\texttt{figures/\_scripts/fig01\_example.py} and \texttt{figures/\_scripts/fig02\_senolytic.py}
(matplotlib). Build: \texttt{make} (or \texttt{tectonic main.tex}); figures: \texttt{make figures}.
The retinal manifest and the senolytic lift are listed verbatim in Appendix~\ref{app:code};
running it reproduces Eq.~\eqref{eq:manifest}, Tab.~\ref{tab:sweep}, and
Figs.~\ref{fig:decomp}--\ref{fig:senolytic} exactly.

\section{Conclusion}
\label{sec:conclusion}
The retinal photoreceptor cure, designed to in-silico exhaustion, reduces to a single binding
constraint --- lost-photoreceptor neogenesis, the largest lost-mass class of any
chronic-degenerative cure --- which caps restoration at $0.64$, below the $\geq 90\%$ gate. The
same senolytic niche-clearance lever that closes the framework's other cures closes the retinal
gate at ${\sim}72\%$ clearance, by lifting lost-class neogenesis and unlocking the M\"uller-glia
reprogramming arm. Uniquely, the predicted lever already exists in the clinic as the local BCL-xL
senolytic UBX1325/foselutoclax. The most consequential next step is therefore not a new target but
a single measurement: does intravitreal senolytic clearance raise human photoreceptor-neogenesis
efficiency toward the ${\sim}0.85$ the gate requires?

\paragraph{Acknowledgments.}
Used the \texttt{hexa-lang} verification surface and the \texttt{sidecar/paper} scaffold.
LLM collaborator: Claude Opus 4.x (Anthropic), which executed the in-silico probes under human
direction; it did not adjudicate correctness.

\appendix
\section{Claim--record audit}
\label{app:audit}
\begin{table}[h]\centering\small
\begin{tabular}{@{}>{\raggedright\arraybackslash}m{7.2cm} >{\raggedright\arraybackslash}m{3.8cm} c@{}}
\toprule
\textbf{Claim} & \textbf{Backing record} & \textbf{Tier} \\
\midrule
generic axis-collapse model & \texttt{cure\_axis\_collapse.py} & \tierGreen \\
retinal ceiling $=0.64$ (gate BLOCK) & \texttt{RETINA-CURE} & \tierGreen \\
binding axis $=$ lost-photoreceptor neogenesis & \texttt{RETINA-CURE} & \tierGreen \\
senolytic $\phi^\ast\!\approx\!72\%$ closes gate & \texttt{RETINA-CURE PD} & \tierGreen \\
per-class $\eta$ + $\phi\!\to\!\eta$ coupling & literature-order & \tierOrange \\
M\"uller$\to$photoreceptor reprogramming arm & \citep{jorstad2017,yao2018} & \tierYellow \\
BCL-xL retinal senolytic node (UBX1325) & \citep{crespodawid2024} & \tierYellow \\
\bottomrule
\end{tabular}
\caption{Every claim maps to a persisted record or citation + tier. The \tierOrange\ row bounds
what is estimated vs established.}
\label{tab:audit}
\end{table}

\paragraph{Falsifier register (pre-registered).} (i) a retinal cure that does NOT block on
lost-photoreceptor neogenesis when reversible+dormant classes are optimised would refute the
binding-axis claim; (ii) senolytic clearance that does NOT raise in-vitro photoreceptor-neogenesis
efficiency would refute the cross-cutting-key claim; (iii) intravitreal BCL-xL senolysis that
cannot reach $\phi^\ast$ retinal clearance, or does so only with systemic platelet toxicity, would
refute the delivery premise.

\section{Generic axis-collapse model + retinal manifest (reproducibility listing)}
\label{app:code}
The entire result is produced by one function applied to the retinal manifest (single dispatch;
no per-disease code). Verbatim core:

\begin{lstlisting}[basicstyle=\ttfamily\footnotesize,frame=single,xleftmargin=0pt,breaklines=true]
def axis_collapse(mass, eta, gate=0.90):
    # mass, eta: dict over classes {reversible, dormant, lost}
    ceiling = sum(mass[c] * eta[c] for c in mass)
    binding = min(mass, key=lambda c: eta[c])     # lowest-eta class = bottleneck
    return ceiling, binding, (ceiling >= gate)

def senolytic_lift(mass, eta, phi, gate=0.90):
    # phi = niche senescent-cell clearance fraction in [0,1]
    e = dict(eta)
    e['lost'] = 0.45 + 0.55*phi          # lost-photoreceptor neogenesis lift
    e['dorm'] = min(0.70 + 0.45*phi, 1)  # Mueller-glia reprogramming arm unlocked
    return axis_collapse(mass, e, gate)

# the retinal instance is a MANIFEST ONLY (no per-disease code):
#   stressed photoreceptors (reversible) | Mueller glia (dorm) | lost photoreceptors (lost)
RETINA = (dict(rev=.25, dorm=.35, lost=.40), dict(rev=.85, dorm=.70, lost=.45))
# axis_collapse(*RETINA) -> ceiling=0.6375, binding='lost', BLOCK
# sweep phi -> gate closes at phi* ~ 0.72
\end{lstlisting}

\noindent Running this reproduces Eq.~\eqref{eq:manifest}, Tab.~\ref{tab:sweep}, and
Figs.~\ref{fig:decomp}--\ref{fig:senolytic}. The full scripts (with the senolytic PD sweep)
are in \texttt{exports/RETINA-CURE/} and \texttt{exports/CURE-PRIMITIVE/round1/}.

\bibliographystyle{plainnat}
\bibliography{references}
\end{document}
